\documentclass[12pt,a4paper]{article}

% ─── Paquetes ───────────────────────────────────────────────────────────────
\usepackage[utf8]{inputenc}
\usepackage[T1]{fontenc}
\usepackage[]{babel}
\usepackage{geometry}
\usepackage{graphicx}
\usepackage{booktabs}
\usepackage{amsmath}
\usepackage{amssymb}
\usepackage{hyperref}
\usepackage{caption}
\usepackage{subcaption}
\usepackage{float}
\usepackage{xcolor}
\usepackage{array}
\usepackage{tabularx}
\usepackage{multirow}
\usepackage{setspace}
\usepackage{parskip}
\usepackage{fancyhdr}
\usepackage{titlesec}

\geometry{
  top=2.5cm, bottom=2.5cm,
  left=2.5cm, right=2.5cm
}

% ─── Encabezado y pie de página ─────────────────────────────────────────────
\pagestyle{fancy}
\fancyhf{}
\rhead{\small Análisis Numérico · UNAL}
\lhead{\small Neural ODEs sobre Datos Reales}
\cfoot{\thepage}
\renewcommand{\headrulewidth}{0.4pt}

% ─── Colores ────────────────────────────────────────────────────────────────
\definecolor{unal}{RGB}{0,60,120}
\definecolor{lightgray}{RGB}{245,245,245}

% ─── Títulos de secciones ───────────────────────────────────────────────────
\titleformat{\section}
  {\large\bfseries\color{unal}}
  {\thesection.}{0.5em}{}[\titlerule]

\titleformat{\subsection}
  {\normalsize\bfseries\color{unal!80}}
  {\thesubsection.}{0.5em}{}

% ─── Hiperlinks ─────────────────────────────────────────────────────────────
\hypersetup{
  colorlinks=true,
  linkcolor=unal,
  urlcolor=unal,
  citecolor=unal
}

% ────────────────────────────────────────────────────────────────────────────
\begin{document}

% ─── Portada ────────────────────────────────────────────────────────────────
\begin{titlepage}
  \centering
  \vspace*{1cm}
  {\Large\bfseries\color{unal} Universidad Nacional de Colombia\par}
  \vspace{0.3cm}
  {\large Análisis Numérico\par}
  {\normalsize Prof.\ Carlos Manuel Orrego Franco\par}
  \vspace{1.5cm}
  \rule{\linewidth}{1.5pt}\\[0.4cm]
  {\Huge\bfseries Neural ODEs sobre Datos Reales\par}
  \vspace{0.3cm}
  {\large COVID-19 JHU --- Nuevos casos diarios (178 países $\times$ 400 días)\par}
  \rule{\linewidth}{1.5pt}
  \vspace{1.5cm}

  \begin{tabular}{ll}
    \textbf{Integrante 1:} & Valentina Aristizábal Sierra \\[4pt]
    \textbf{Integrante 2:} & Camilo Andrés Jaime Rojas \\[4pt]
    \textbf{Entrega:}      & 27 de mayo de 2026 \\
  \end{tabular}

  \vfill
  {\small Actividad Final --- Parejas}
\end{titlepage}

% ─── Tabla de contenidos ────────────────────────────────────────────────────
\tableofcontents
\newpage

% ────────────────────────────────────────────────────────────────────────────
\section{Descripción del dataset}
\label{sec:dataset}

El dataset utilizado proviene del repositorio COVID-19 del Center for Systems
Science and Engineering (CSSE) de la Universidad Johns Hopkins (JHU), que
consolidó diariamente casos confirmados a nivel mundial desde enero de 2020.
Para este experimento se procesó la serie de nuevos casos diarios para 178
países, abarcando 400 días consecutivos.
La variable es unidimensional ($\text{dim}=1$), correspondiente a nuevos casos
normalizados mediante transformación logarítmica suavizada, con rango efectivo
$[-1.098,\;7.928]$.
El conjunto se dividió en 143 trayectorias de entrenamiento y 35 de prueba.

Desde el punto de vista dinámico, las series de COVID-19 exhiben el comportamiento
oscilatorio y no lineal que motiva el uso de EDOs para su modelamiento.
Al igual que el modelo SIR estudiado en los notebooks NB03--NB04 del curso
($dI/dt = \beta S I - \gamma I$), la dinámica de contagio presenta crecimiento
exponencial inicial, un pico y una caída.
Sin embargo, a diferencia del SIR clásico, las trayectorias reales muestran
múltiples oleadas, efectos de política pública (confinamientos, vacunación)
y heterogeneidades profundas entre países --- factores no observables en la
serie de casos, como se aprecia en la Figura~\ref{fig:exploracion}.

Esta variabilidad representa el principal desafío del dataset: países con
dinámicas muy distintas (olas rápidas y agudas versus comportamientos prolongados
multi-pico) coexisten en el mismo conjunto.
El muestreo es formalmente regular (400 pasos igualmente espaciados), pero el
ruido inherente a los reportes epidemiológicos (subregistro, cambios en criterios
diagnósticos, retardo en notificación) introduce una irregularidad implícita que
pone a prueba la capacidad de los modelos de aprender señal sobre ruido.

\begin{figure}[H]
  \centering
  \includegraphics[width=\linewidth]{fig_01_exploracion.png}
  \caption{Exploración del dataset: 8 trayectorias de países representativos
    (nuevos casos diarios normalizados). Se observa la diversidad de patrones ---
    olas únicas pronunciadas, picos múltiples y señales planas --- que refleja
    la heterogeneidad epidemiológica global.}
  \label{fig:exploracion}
\end{figure}

% ────────────────────────────────────────────────────────────────────────────
\section{Hiperparámetros elegidos y justificación}
\label{sec:hiperparametros}

Se ejecutaron dos configuraciones para explorar el efecto del tamaño del modelo.
La \textbf{Configuración A} ($\texttt{dim\_latente}=8$) apuesta por mayor
capacidad representacional: el vector $z_0$ puede codificar simultáneamente
información sobre la fase de la ola, la tasa de crecimiento y la
heterogeneidad entre países.
La \textbf{Configuración B} ($\texttt{dim\_latente}=4$) opera con la premisa
de que un espacio latente más pequeño funciona como regularizador implícito,
forzando al modelo a aprender solo las componentes más explicativas de la
dinámica --- relevante dado el alto nivel de ruido del dataset.

\begin{table}[H]
  \centering
  \caption{Hiperparámetros de ambas configuraciones y justificación.}
  \label{tab:hiperparametros}
  \small
  \begin{tabularx}{\linewidth}{lccX}
    \toprule
    \textbf{Hiperparámetro} & \textbf{Config A} & \textbf{Config B} & \textbf{Justificación} \\
    \midrule
    \texttt{dim\_latente}    & 8    & 4    & Dimensión de $z_0$. Config A captura más modos de variabilidad; Config B regulariza para datos ruidosos. \\[4pt]
    \texttt{n\_pasos\_ode}   & 10   & 6    & Pasos RK4. Más pasos $\Rightarrow$ mayor precisión numérica pero mayor costo $\mathcal{O}(n)$. \\[4pt]
    \texttt{learning\_rate}  & 0.001 & 0.001 & Tasa Adam estándar para Neural ODEs \cite{chen2018}. Tasas mayores producen inestabilidades en el solver RK4. \\[4pt]
    \texttt{n\_epochs}       & 35   & 35   & La función de costo se estabilizó $\approx$ época 30 en ambas configs (ver Figura~\ref{fig:curvas}). \\[4pt]
    \texttt{batch\_size}     & 32   & 32   & Balance entre varianza del gradiente y memoria de GPU. Estándar en la literatura. \\[4pt]
    \texttt{dim\_oculto\_f}  & 48   & 32   & Capacidad de $f_\theta$. Proporcional a \texttt{dim\_latente} para mantener expresividad. \\[4pt]
    \texttt{dim\_oculto\_enc}& 32   & 24   & Unidades del encoder RNN inverso. Coherente con $z_0$. \\[4pt]
    \texttt{dim\_oculto\_dec}& 48   & 32   & Capacidad del decoder; $\geq$ \texttt{dim\_oculto\_f} para reconstrucción fiel. \\
    \bottomrule
  \end{tabularx}
\end{table}

% ────────────────────────────────────────────────────────────────────────────
\section{Tabla de resultados (RMSE)}
\label{sec:resultados}

Las tablas siguientes reportan el RMSE de reconstrucción (\emph{rec}) y
extrapolación (\emph{ext}) para cada modelo y fracción de observaciones,
replicando el esquema experimental de la Tabla~2 de \citet{rubanova2019}.
El valor ``---'' en extrapolación al 100\% es esperado: con todos los puntos
observados no quedan puntos futuros para evaluar.

\begin{table}[H]
  \centering
  \caption{Configuración A: \texttt{dim\_latente}$=8$, \texttt{n\_pasos\_ode}$=10$,
           \texttt{dim\_oculto\_f}$=48$.}
  \label{tab:config_a}
  \begin{tabular}{lcccccc}
    \toprule
    \multirow{2}{*}{\textbf{Modelo}}
      & \multicolumn{2}{c}{\textbf{30\% obs}}
      & \multicolumn{2}{c}{\textbf{60\% obs}}
      & \textbf{100\% obs}
      & \\
    \cmidrule(lr){2-3}\cmidrule(lr){4-5}\cmidrule(lr){6-6}
      & rec & ext & rec & ext & rec & \textbf{Paráms} \\
    \midrule
    Método Clásico (splines) & 0.000 & 0.114 & 0.000 & 0.087 & 0.000 & 4 \\
    Neural ODE               & 0.500 & 0.501 & 0.498 & 0.505 & 0.501 & 3\,721 \\
    Latent ODE               & 0.576 & 0.576 & 0.573 & 0.579 & 0.576 & 7\,593 \\
    \bottomrule
  \end{tabular}
\end{table}

\begin{table}[H]
  \centering
  \caption{Configuración B: \texttt{dim\_latente}$=4$, \texttt{n\_pasos\_ode}$=6$,
           \texttt{dim\_oculto\_f}$=32$.}
  \label{tab:config_b}
  \begin{tabular}{lcccccc}
    \toprule
    \multirow{2}{*}{\textbf{Modelo}}
      & \multicolumn{2}{c}{\textbf{30\% obs}}
      & \multicolumn{2}{c}{\textbf{60\% obs}}
      & \textbf{100\% obs}
      & \\
    \cmidrule(lr){2-3}\cmidrule(lr){4-5}\cmidrule(lr){6-6}
      & rec & ext & rec & ext & rec & \textbf{Paráms} \\
    \midrule
    Método Clásico (splines) & 0.000 & 0.114 & 0.000 & 0.087 & 0.000 & 4 \\
    Neural ODE               & 0.503 & 0.504 & 0.501 & 0.508 & 0.504 & 1\,581 \\
    Latent ODE               & 0.561 & 0.561 & 0.559 & 0.565 & 0.561 & 3\,717 \\
    \bottomrule
  \end{tabular}
\end{table}

El Método Clásico (splines) logra RMSE de reconstrucción exactamente 0.000 en todos
los casos --- los splines interpolantes pasan exactamente por los puntos observados.
Sin embargo, su RMSE de extrapolación es considerablemente más alto
($\approx 0.114$ con 30\% de observaciones).
La Neural ODE mantiene un comportamiento consistente entre reconstrucción y
extrapolación ($\sim$0.50), lo que indica que no memoriza sino que aprende una
dinámica general.
El Latent ODE presenta valores ligeramente superiores a la Neural ODE en ambas
configuraciones.

% ────────────────────────────────────────────────────────────────────────────
\section{Figuras}
\label{sec:figuras}

\subsection{Comparación de modelos (60\% de observaciones)}

\begin{figure}[H]
  \centering
  % \includegraphics[width=\linewidth]{figuras/fig2_comparacion.png}
  \includegraphics[width=\linewidth]{fig_02_comparacion.png}
  \caption{Comparación de los tres modelos con 60\% de observaciones en
    3 trayectorias de prueba.
    Los puntos son las observaciones disponibles, la línea discontinua es la
    predicción y la línea sólida gris es la trayectoria real.
    Fila superior: Método Clásico (splines).
    Fila media: Neural ODE.
    Fila inferior: Latent ODE.}
  \label{fig:comparacion}
\end{figure}

\subsection{Curvas de aprendizaje}

\begin{figure}[H]
  \centering
  \includegraphics[width=0.85\linewidth]{fig_03_curvas.png}
  \caption{Curvas de aprendizaje para Neural ODE (verde) y Latent ODE (azul)
    durante 35 épocas.
    La Neural ODE converge de forma más suave ($\text{loss}_0\approx 0.35$);
    el Latent ODE muestra una caída inicial pronunciada desde
    $\text{loss}_0\approx 0.55$ seguida de estabilización $\approx$ época~30.}
  \label{fig:curvas}
\end{figure}

% ────────────────────────────────────────────────────────────────────────────
\section{Respuestas a las preguntas de análisis}
\label{sec:preguntas}

\subsection{P1 — ¿Con qué porcentaje de observaciones empieza a ganar la
  Neural ODE o el Latent ODE?}
\label{sec:p1}

Con base en los resultados obtenidos (Tablas~\ref{tab:config_a}
y~\ref{tab:config_b}), la Neural ODE \textbf{no supera al Método Clásico en
extrapolación} en ninguna de las fracciones evaluadas (30\%, 60\%, 100\%) para
este dataset.
Con 30\% de observaciones, el Método Clásico registra RMSE de extrapolación de
0.114 frente a $\sim$0.501--0.504 de la Neural ODE.
Con 60\%, el clásico mejora a 0.087 mientras la Neural ODE se mantiene
alrededor de 0.505--0.508.

Este resultado es coherente con lo trabajado en los NB06--07.
Los splines cúbicos son interpoladores exactos y, para series epidémicas con
continuidad local, su extrapolación a corto plazo puede ser precisa.
La Neural ODE aprende una dinámica \emph{promedio} sobre 143 países de
entrenamiento y, al predecir en los 35 países de prueba, esa dinámica
generalizada puede no coincidir con el régimen local de cada país.
En otras palabras, el modelo clásico tiene la ventaja de ajustarse
específicamente a cada trayectoria, mientras la Neural ODE generaliza desde el
entrenamiento; en datasets con dinámica muy heterogénea entre trayectorias
(como COVID-19 global), esa generalización penaliza el rendimiento individual.
La ventaja de la Neural ODE se manifestaría con más épocas, mayor capacidad, o
en datasets donde la dinámica compartida entre trayectorias es más fuerte
(como el SIR sintético de la Opción A).

\subsection{P2 — ¿Por qué el encoder RNN corre en orden inverso al tiempo?}
\label{sec:p2}

Esta es una de las decisiones de diseño más elegantes del Latent ODE
\cite{rubanova2019}.
El encoder RNN procesa la secuencia $\{x_T, x_{T-1}, \ldots, x_1\}$ de atrás
hacia adelante para poder computar el estado inicial $z_0$ en el tiempo $t_0$
como el estado final del procesamiento.
La lógica es directa: el Latent ODE modela la trayectoria como la solución de
una EDO integrada desde $z_0$ hacia adelante en el tiempo.
Para inferir $z_0$, el encoder necesita condensar toda la información observada
en un vector de condición inicial, y $t_0$ es el punto más antiguo de la
secuencia.

Si el encoder corriese en orden cronológico normal, su estado oculto al
finalizar correspondería al tiempo más reciente ($t_T$), no al más antiguo ---
lo que haría incoherente la integración hacia adelante desde $t_0$.
Además, por el fenómeno de \emph{olvido selectivo} del RNN (estudiado en NB05),
el estado final estaría dominado por las observaciones más recientes, perdiendo
información sobre las condiciones iniciales de la dinámica.
En COVID-19, esto significaría que $z_0$ reflejaría el estado actual de la ola
pero olvidaría la pendiente inicial de crecimiento --- precisamente la
información necesaria para integrar la EDO desde el inicio.

\subsection{P3 — ¿Hubo algún resultado sorpresivo?}
\label{sec:p3}

El resultado más llamativo es que el \textbf{Latent ODE no supera a la Neural
ODE} sino que obtiene RMSE consistentemente más alto en todas las condiciones
($\sim$0.576 vs.\ $\sim$0.501 en Config~A; $\sim$0.561 vs.\ $\sim$0.503 en
Config~B).
Intuitivamente se esperaría lo contrario: el Latent ODE es un modelo más
sofisticado con encoder RNN y espacio latente probabilístico, diseñado para
capturar la heterogeneidad entre trayectorias y producir predicciones
personalizadas.

Hay varias razones que explican este resultado:
\begin{enumerate}
  \item La curva de aprendizaje (Figura~\ref{fig:curvas}) muestra que el Latent
    ODE parte de una pérdida inicial considerablemente más alta
    ($\sim$0.55 vs.\ $\sim$0.35), lo que con solo 35 épocas resulta en
    convergencia más lenta e incompleta.
  \item El dataset COVID-19 tiene una única dimensión y muestreo regular, lo
    que reduce la ventaja del encoder RNN inverso: al no haber irregularidad
    temporal real, la información que el encoder debe condensar es más simple
    y la Neural ODE puede aproximarla adecuadamente.
  \item El ruido epidemiológico dificulta que el encoder aprenda un $z_0$
    representativo, degradando toda la integración subsecuente.
\end{enumerate}

Este resultado es consistente con \citet{rubanova2019}: la ventaja del Latent
ODE se manifiesta principalmente con muestreo genuinamente irregular y variables
multidimensionales, como en PhysioNet.

\subsection{P4 — ¿En qué tipo de problema NO recomendaríamos una Neural ODE?}
\label{sec:p4}

Con base en los resultados de este experimento, identificamos tres escenarios en
los que una Neural ODE no es la herramienta adecuada:

\begin{enumerate}
  \item \textbf{Cuando los datos son ruidosos y la interpolación exacta es
    suficiente.}
    El Método Clásico obtuvo RMSE de reconstrucción exactamente 0 en todos los
    casos frente al $\sim$0.50 de la Neural ODE, siendo prácticamente instantáneo
    (0.0~s vs.\ 2\,473~s).
    Si el objetivo es rellenar valores faltantes en una serie donde los puntos
    observados son confiables, los splines son superiores sin ningún costo
    computacional.

  \item \textbf{Cuando el presupuesto computacional es limitado.}
    Con 143 trayectorias de entrenamiento, la Neural ODE requirió
    $\sim$41~minutos y el Latent ODE $\sim$54~minutos en Colab.
    En un contexto de epidemiología operativa donde se necesitan actualizaciones
    diarias, este costo es inviable.
    El número de parámetros (3\,721--7\,593) relativamente alto respecto al
    tamaño del dataset también puede inducir sobreajuste sin regularización
    explícita.

  \item \textbf{Cuando la dinámica entre trayectorias es extremadamente
    heterogénea sin un campo vectorial $f_\theta$ común.}
    En COVID-19 global, las diferencias entre países responden a factores
    externos no observables en la serie de casos.
    Una Neural ODE que aprende una $f_\theta$ compartida promedia estas
    diferencias.
    En este caso, métodos que ajustan parámetros por trayectoria (como el SIR
    con \texttt{scipy.optimize}) pueden ser más precisos porque se especializan
    en cada contexto individual.
\end{enumerate}

% ────────────────────────────────────────────────────────────────────────────
\section{Conclusiones}
\label{sec:conclusiones}

Esta actividad permitió materializar el recorrido completo desde la teoría de
EDOs hasta la aplicación de Neural ODEs y Latent ODEs sobre datos
epidemiológicos reales.
Los resultados muestran que, en el dataset COVID-19 con muestreo regular y una
sola dimensión, el método clásico con splines sigue siendo competitivo en
extrapolación a corto plazo, mientras que la Neural ODE logra una dinámica más
robusta y consistente entre reconstrucción y extrapolación, sin memorizar.

El Latent ODE, a pesar de su mayor sofisticación, no superó a la Neural ODE
aquí --- probablemente por convergencia más lenta con pocas épocas y por la
ausencia de irregularidad temporal real, que es precisamente el escenario en
que brilla según \citet{rubanova2019}.

La comparación de dos configuraciones confirmó que reducir el espacio latente
(Config~B) ofrece un compromiso razonable: resultados casi equivalentes con la
mitad de parámetros y menor tiempo de entrenamiento.
En síntesis, las Neural ODEs son herramientas poderosas cuando la dinámica es
genuinamente compartida entre trayectorias y el muestreo es irregular ---
condiciones que posicionan al Latent ODE como el enfoque de referencia para
series fisiológicas multivariadas como las del dataset PhysioNet, el escenario
canónico del paper de \citet{rubanova2019}.

% ────────────────────────────────────────────────────────────────────────────
\section*{Referencias}
\addcontentsline{toc}{section}{Referencias}

\begin{thebibliography}{9}

\bibitem{chen2018}
  Chen, R.~T.~Q., Rubanova, Y., Bettencourt, J., \& Duvenaud, D. (2018).
  Neural Ordinary Differential Equations.
  \textit{Advances in Neural Information Processing Systems (NeurIPS)}, 31.

\bibitem{rubanova2019}
  Rubanova, Y., Chen, R.~T.~Q., \& Duvenaud, D. (2019).
  Latent ODEs for Irregularly-Sampled Time Series.
  \textit{Advances in Neural Information Processing Systems (NeurIPS)}, 32.

\bibitem{dong2020}
  Dong, E., Du, H., \& Gardner, L. (2020).
  An interactive web-based dashboard to track COVID-19 in real time.
  \textit{The Lancet Infectious Diseases}, 20(5), 533--534.

\end{thebibliography}

\end{document}
